\section*{\SGAInlineRefTarget{sga3:x:section:7}7. 用扩大基本群分类扭曲常预概形与有限型乘法型群}
\addcontentsline{toc}{section}{7. 用扩大基本群分类}

\setcounter{footnote}{57}

\subsection*{\SGAInlineRefTarget{sga3:x:subsection:7-0}7.0.}

设 \(S\) 是预概形，仍假定它局部 Noether，以保证 \(S\) 及我们要考虑的
\(S\) 上某些预概形（特别是在 \(S\) 上平展的预概形，更一般地在 \(S\) 上局部有限型的预概形）局部连通。

\par\medskip\noindent\SGAInlineRefTarget{sga3:x:proposition:7-0-1}\textbf{命题 7.0.1。---}\enspace
\emph{\(S\) 上对于 \((\mathrm{fppf})\) 拓扑局部平凡的每个扭曲常预概形 \(X\)
（即由一个忠实平坦且局部有限呈示的态射 \(S'\to S\) 分裂的预概形）都是拟等平凡的
（即甚至可选择 \(S'\to S\) 为平展满射）。}

事实上，可假定 \(S\) 连通，从而 \(X\) 具有型 \(I\)，其中 \(I\) 是固定的集合。因此
\[
  X'=X\times_S S'
\]
同构于 \(I_{S'}\)，故 \(I_{S'}\) 配备关于 \(S'\to S\) 的下降数据；换言之，有同构
\[
  I_{S''}\xrightarrow{\sim}I_{S''}
\]
满足通常的传递性条件。现在 \(S''=S'\times_S S'\) 局部 Noether，从而局部连通。
因此，\(I_{S''}\) 的自同构对应于 \(G_{S''}\) 的截面，其中
\(G=\operatorname{Aut}(I)\) 是 \(I\) 的置换群。

由此在 \(G_{S'}\)（视为平凡 Galois 主丛）上得到关于 \(S'\to S\) 的下降数据。
由 \SGARefLink{sga3:x:lemma:5-4}{5.4}，这些下降数据有效，因而给出 \(S\) 上群为
\(G\) 的 Galois 主丛 \(P\)。由构造，\(P\) 在 \(S\) 上局部 Noether 的预概形范畴中
表示函子 \(\operatorname{Isom}_S(I_S,X)\)。因此，平展满射基变换 \(P\to S\)
分裂 \(X\)，所以 \(X\) 确为拟等平凡的。

\par\medskip\noindent\SGAInlineRefTarget{sga3:x:remark:7-0-2}\textbf{注记 7.0.2。---}\enspace
应牢记，即使 \(S\) 是域的谱，一般也不能断言 \(S\) 上的每个扭曲常丛都拟等平凡。
例如，可取 \(X\) 为一列形如 \(\operatorname{Spec}(k_i)\) 的概形的不交并，
其中各 \(k_i\) 是 \(k\) 的次数严格递增的可分扩张。

刚才的证明同时表明，\(S\) 上型为 \(I\) 的拟等平凡扭曲常丛 \(X\) 的分类，
等价于 \(S\) 上群为 \(G=\operatorname{Aut}(I)\) 的 Galois 主丛的分类。
事实上，这是范畴的等价。

如同 SGA~1，V，可将它表述为更方便的形式。为此，假定 \(S\) 连通，且配备几何点
\(\xi\)。于是定义了扩大基本预群 \(\Pi=\Pi_1(S,\xi)\)。对 \(S\) 上的每个拟等平凡
扭曲常丛 \(X\)，令 \(I=X(\xi)\) 为其在 \(\xi\) 处的集合论纤维。
则 \(X\) 具有型 \(I\)，因而如上与 \(S\) 上群为 \(G=\operatorname{Aut}(I)\)
的 Galois 主丛相伴：
\[
  P=\operatorname{Isom}_S(I_S,X)
\]
。

由 \(\Pi\) 的定义，因此得到从 \(\Pi\) 到 \(G\) 的典范同态，即从某个
\(\Pi_i\) 到 \(G\) 的同态。由于 \(G\) 是 \(I=X(\xi)\) 的置换群，
这表示 \(\Pi\)“连续作用于 \(I=X(\xi)\)”，即当 \(i\) 充分大时，
各 \(\pi_i\) 在 \(I\) 上的作用与过渡态射相容。

留待读者验证，\(S\) 上拟等平凡扭曲常丛之间的每个 \(S\)-态射 \(X\to Y\)
都诱导一个与 \(\Pi\) 的作用相容的映射 \(X(\xi)\to Y(\xi)\)，
而由此得到的函子是范畴等价：

\par\medskip\noindent\SGAInlineRefTarget{sga3:x:proposition:7-0-3}\textbf{命题 7.0.3。---}\enspace
\emph{设 \(S\) 是局部 Noether 连通预概形，\(\xi\) 是 \(S\) 的几何点，
\(\Pi=\Pi_1(S,\xi)\) 是 \(S\) 在 \(\xi\) 处的扩大基本预群。则函子}
\[
  X\longmapsto X(\xi)
\]
\emph{是 \(S\) 上拟等平凡扭曲常丛范畴与 \(\Pi\) 连续作用的集合范畴之间的等价。}

此函子与有限和及有限逆极限相容。例如，\(S\) 上的拟等平凡扭曲常群
（或环等等）对应于预群 \(\Pi\) 连续作用的通常群（分别为环等等）。特别地：

\par\medskip\noindent\SGAInlineRefTarget{sga3:x:corollary:7-0-4}\textbf{推论 7.0.4。---}\enspace
\emph{\(S\) 上的拟等平凡扭曲常交换群范畴等价于“\(\Pi\)-模”范畴，
即 \(\Pi\) 连续作用的交换群 \(M\) 的范畴。}

使用 \SGARefLink{sga3:x:corollary:5-7}{5.7}，得到：

\par\medskip\noindent\SGAInlineRefTarget{sga3:x:theorem:7-1}\textbf{定理 7.1。---}\enspace
\emph{设 \(S\) 是局部 Noether 连通预概形，\(\xi\) 是 \(S\) 的几何点，
\(\Pi=\Pi_1(S,\xi)\) 是 \(S\) 在 \(\xi\) 处的扩大基本预群。则函子}
\[
  G\longmapsto
  \operatorname{Hom}_{\kappa(\xi)\text{-}\mathrm{gr}}
    (G_\xi,\mathbf{G}_{m,\xi})
\]
\emph{在 \(S\) 上的拟等平凡乘法型群范畴与 \(\Pi\)-模范畴之间诱导反等价。}

使用 \SGARefLink{sga3:x:corollary:4-5}{4.5}，得到：

\par\medskip\noindent\SGAInlineRefTarget{sga3:x:corollary:7-2}\textbf{推论 7.2。---}\enspace
\emph{前述函子在 \(S\) 上的有限型乘法型群范畴与在 \(\mathbb{Z}\) 上有限型的
\(\Pi\)-模范畴之间诱导反等价。}

\par\medskip\noindent\SGAInlineRefTarget{sga3:x:example:7-3}\textbf{例 7.3。---}\enspace
例如，将 \(S\) 取为代数闭域上的完备有理曲线，恰有一个具有 \(n+1\) 个不同分支的多重点。
由 \SGARefLink{sga3:x:example:6-4}{6.4}，扩大基本群 \(\Pi(S,\xi)\) 是 \(n\) 个生成元上的自由群。
因此由 \SGARefLink{sga3:x:corollary:7-2}{7.2}，\(S\) 上相对维数为 \(m\) 的环面的分类，
等价于 \(\mathbb{Z}\)-模 \(\mathbb{Z}^m\) 的 \(n\) 个自同态 \(A_1,\ldots,A_n\)
构成的系在 \(\mathbb{Z}^m\) 的自同构意义下的分类。除非 \(n\leq 1\) 或 \(m\leq 1\)，
这种系的显式分类似乎无望。至少可以为此类系定义大量非平凡不变量，例如各 \(A_i\)
的特征多项式。\footnote{\textit{编者注。}这特别表明
\SGARefLink{sga3:x:paragraph:1:6}{1.6} 中考虑的相对维数为 \(2\) 的环面非等平凡。}

\par\medskip\noindent\SGAInlineRefTarget{sga3:x:remark:7-4}\textbf{注记 7.4。---}\enspace
对 \(S\) 不作任何假设，仍有如下事实：对于给定的在 \(\mathbb{Z}\) 上有限型的通常交换群
\(M\)，\(S\) 上型为 \(M\) 的乘法型群范畴反等价于 \(S\) 上群为
\(G=\operatorname{Aut}_{\mathrm{gr}}(M)\) 的 Galois 主丛范畴。
这由 \SGARefLink{sga3:x:corollary:5-9}{5.9} 和 \SGARefLink{sga3:x:corollary:5-10}{5.10} 容易推出。

\par\medskip\noindent\SGAInlineRefTarget{sga3:x:remark:7-5}\textbf{注记 7.5。---}\enspace
这两节中简述的基本预群理论，更适于在一般景的框架中表述。
这样，它同样适用于例如通常的拓扑空间，至少对于局部连通拓扑空间
（不必局部单连通）给出令人满意的理论。这里似乎也不能只使用基本群，而需要一个预群。
最后，一旦掌握了拓扑与下降的语言（它确实处于这些问题的核心），
这里简述的论述在技术上也比 SGA~1，V 更简单，因此原则上应取代后者。

\section*{\SGAInlineRefTarget{sga3:x:section:8}8. 附录：去掉某些仿射性假设}
\addcontentsline{toc}{section}{8. 附录：去掉某些仿射性假设}

目的是证明 4.9 的下列推广。

\par\medskip\noindent\SGAInlineRefTarget{sga3:x:theorem:8-1}\textbf{定理 8.1。---}\enspace
\emph{设 \(S\) 是预概形，\(H\) 是平坦有限呈示的 \(S\)-群预概形，
具有连通仿射纤维。\footnote{\textit{编者注。}注意，由 VI\textsubscript{B}，
\SGARefLink{thm:sga3-VIB-5.3}{定理~5.3} 或 \SGARefLink{sga3:x:corollary:5-5}{推论~5.5}，
这些假设蕴含 \(H\) 在 \(S\) 上分离。}设 \(s\in S\)，假定 \(H_s\) 为环面。
则存在 \(s\) 的开邻域 \(U\)，使 \(H|_U\) 为环面。}

立即得到：

\par\medskip\noindent\SGAInlineRefTarget{sga3:x:corollary:8-2}\textbf{推论 8.2。---}\enspace
\emph{设 \(H\) 是在 \(S\) 上平坦且有限呈示的 \(S\)-群预概形。
则 \(H\) 为环面，当且仅当其纤维为环面。}

\par\medskip\noindent\SGAInlineRefTarget{sga3:x:remark:8-3}\textbf{注记 8.3。---}\enspace
即使 \(S\) 是离散赋值环的谱，也不能去掉 8.1 中关于 \(H\) 的纤维
（此处为泛纤维）仿射的假设：存在 \(S\) 上的平滑群，其泛纤维为椭圆曲线，
特殊纤维为 \(\mathbf{G}_m\)。

\medskip
\noindent\emph{8.1 的证明。}
显然可假定 \(S=\operatorname{Spec}(A)\) 仿射。
由标准步骤（参见\ EGA~IV\textsubscript{2}，8.8.2），这将问题归约到 \(S\) 还是
Noether 概形的情形。先在此情形中证明 8.2。

由 4.9，只须证明 \(H\) 在 \(S\) 上仿射。因此可以假定
\footnote{\textit{编者注。}再次使用 EGA~IV\textsubscript{2}，8.8.2。}
\(A\) 是 Noether 局部环。由于其完备化 \(\widehat A\) 在 \(A\) 上忠实平坦，
通过下降归约到 \(A\) 为完备 Noether 局部环的情形。
设 \(S'\) 是 \(S_{\mathrm{red}}\) 的正规化。Nagata 定理断言它在 \(S\) 上有限
（EGA~\(0_{\mathrm{IV}}\)，23.1.5）。此外，\(S'\to S\) 为满射，所以
\[
  H'=H\times_S S'\longrightarrow H
\]
有限且满射。因此，为证明 \(H\) 仿射，只须说明 \(H'\) 仿射（EGA~II，6.7.1）。

将 \(S\) 换成 \(S'\) 的一个连通分量，归约到 \(A\) 为 Noether（半局部）、
正规整环的情形。\footnote{\textit{编者注。}下文扩充了原文。}
此外，必要时将 \(A\) 换成其在分式域的某个有限可分扩张中的正规化后，
可假定 \(H\) 的泛纤维 \(H_\eta\) 可对角化，即存在同构
\[
  u_\eta:H_\eta\xrightarrow{\sim}T_\eta,
  \qquad T=\mathbf{G}_{m,S}^{\,r}.
\]
现在有：

\par\medskip\noindent\SGAInlineRefTarget{sga3:x:lemma:8-4}\textbf{引理 8.4。---}\enspace
\emph{设 \(S\) 是局部 Noether 正规不可约预概形，泛点为 \(\eta\)。
设 \(H\) 是 \(S\) 上的平滑群预概形，具有连通纤维，\(T\) 是 \(S\) 上的
有限型乘法型群预概形，\(u_\eta:H_\eta\to T_\eta\) 是 \(\kappa(\eta)\)
上的代数群同态。则 \(u_\eta\) 延拓为群同态 \(u:H\to T\)。}

必要时将 \(S\) 换成其在函数域的某个有限可分扩张中的正规化后，
可假定 \(T\) 可对角化（在拟作的应用中已经如此）。于是 \(T\) 是形如
\(\mathbf{G}_{m,S}^{\,r}\) 的群的闭子群，从而归约到 \(T=\mathbf{G}_{m,S}\) 的情形。

只须证明，将 \(u_\eta\) 看作从 \(H\) 到 \(T=\mathbf{G}_{m,S}\) 的有理映射时，
它处处有定义（于是延拓它的态射 \(u:H\to T\) 必为群同态）。
可将 \(u_\eta\) 看作 \(H_\eta\) 的结构层的可逆截面 \(f\)，须证明它延拓为
\(H\) 的结构层的可逆截面。由于 \(H\) 在正规预概形 \(S\) 上平滑，
它正规（SGA~1，II 3.1）。因此，只须找到 \(H\) 的余维至少为 \(2\) 的闭子集
\(Z\)，使 \(f\) 延拓为 \(H-Z\) 的结构层的可逆截面。
在 \(S\) 的余维一点处局部化，立即归约到 \(S\) 为离散赋值环 \(A\) 的谱的情形。

设 \(t\) 是 \(A\) 的素元，\(t'\) 是它所定义的 \(\mathcal O_H\) 的截面，
从而特殊纤维 \(H_0\) 等于 \(V(t')\)。由假设，\(H_0\) 在剩余域 \(k\)
上平滑且连通。函数 \(f\) 是 \(H\) 上在 \(H-H_0\) 中无零点也无极点的有理函数。
由于 \(H_0=V(t')\) 是不可约除子，存在整数 \(n\in\mathbb{Z}\)，使
\[
  t'^{\,n}f=t^n f
\]
无零点也无极点，因而是 \(\mathcal O_H\) 的可逆截面。
所以它定义态射 \(v:H\to\mathbf{G}_{m,S}\)。由于 \(v_\eta=t'^{\,n}u_\eta\)，
且 \(u_\eta\) 为群同态，\(v\) 将 \(H\) 的单位截面映为 \(\mathbf{G}_{m,S}\)
的一个截面，其在 \(S\) 的泛点处的值为 \(t^n\)。由于这是 \(\mathbf{G}_{m,S}\)
的截面，\(t^n\) 必为单位；故 \(n=0\)。所以 \(v\) 延拓 \(u_\eta\)，
完成了 \SGARefLink{sga3:x:lemma:8-4}{8.4} 的证明。

将此引理应用于当前情形，得到群同态
\[
  u:H\longrightarrow T=\mathbf{G}_{m,S}^{\,r}
\]
它在泛纤维上诱导同构。将证明 \(u\) 是同构。

\par\medskip\noindent\SGAInlineRefTarget{sga3:x:lemma:8-5}\textbf{引理 8.5。---}\enspace
\emph{对与 \(A\) 的剩余特征互素的每个整数 \(n>0\)，群
\[
  {}_nH=\ker(n\cdot\operatorname{id}_H)
\]
在 \(S\) 上有限。}

若 \(n\) 与 \(A\) 的剩余特征互素，它也与 \(S\) 的各点处的所有剩余特征互素。
因此，\(n\cdot\operatorname{id}_H\) 在 \(H\) 的每个纤维上诱导平展态射，
从而本身平展（SGA~1，I 5.9）；所以其核 \({}_nH\) 在 \(S\) 上平展。
另一方面，\({}_nH\) 在 \(S\) 上分离，因为 \(H\) 如此。
\begingroup
\renewcommand{\thefootnote}{\(\ast\)}
\footnote{由 Raynaud 定理，VI\textsubscript{B}，5.3。}
\endgroup
\addtocounter{footnote}{-1}%
\footnote{\textit{编者注。}应结合 VI\textsubscript{B}，\S5 的修改展开这一点。}
此外，它的所有纤维都有相同的秩 \(n^r\)，因为 \(H\) 的纤维都是维数相同的
\(r\) 维环面（由于 \(H\) 在 \(S\) 上平滑）。因此 \({}_nH\) 在 \(S\) 上有限
（SGA~1，I 10.9，或 EGA~IV\textsubscript{4}，18.2.9）。

由前面的引理，\(u({}_nH)\) 是 \(T\) 的闭子集。在泛纤维 \(T_\eta\) 上，
此子集与 \({}_n(T_\eta)\) 相同，因此包含后者的闭包，即 \({}_nT\)。
对每个 \(s\in S\)，群 \({}_n(T_s)\) 随 \(n\) 变化在 \(T_s\) 中稠密；
由于 \(u_s(H_s)\) 是包含它们的闭子集（VI\textsubscript{B}，1.4.2），
有 \(u_s(H_s)=T_s\)。所以 \(u\) 满射。域上同维环面之间的满同态是平坦的，
\footnote{\textit{编者注。}例如，这由 VIII 3.2(a) 推出。更一般地，若 \(f:G\to H\)
是域 \(k\) 上代数群之间的满射态射，且 \(H\) 约化，则 \(f\) 平坦
（参见\ VI\textsubscript{B}，1.3）。}故 \(u\) 在每个纤维上诱导平坦态射。
因此 \(u\) 平坦（SGA~1，I 5.9）。从而 \(K=\ker(u)\) 在 \(S\) 上平坦，
\footnote{\textit{编者注。}下文扩充了原文。}因而等于其泛纤维 \(K_\eta\) 的闭包。
由构造，\(K_\eta\) 是单位群。由于 \(K\) 在 \(S\) 上分离（因为 \(H\) 如此），
它的单位截面闭，所以 \(K\) 是单位群。因此 \(u\) 是单态射；由于已见它忠实平坦，
它是同构（参见\ SGA~1，I 5.1，或 EGA~IV\textsubscript{4}，17.9.1）。
这证明 \(H\) 是环面，并完成了 8.2 的证明。\footnote{\textit{编者注。}
至少在至此所考虑的情形，即 \(S\) 局部 Noether 时如此。}

\par\medskip\noindent\SGAInlineRefTarget{sga3:x:remark:8-5-1}\textbf{注记 8.5.1。---}\enspace
也可以不引用 8.5，而引用 Zariski“主定理”，它直接蕴含 \(u\) 是开浸入，
从而是同构。\footnote{\textit{编者注。}编者未能理解此注记，
因为看不出为何 \(u\) \emph{先验地}具有有限且满射的纤维。}

为证明 \SGARefLink{sga3:x:theorem:8-1}{8.1}，拟等平凡性定理将问题归约到
\(S\) 局部、\(s\) 为其闭点的情形，\footnote{\textit{编者注。}展开此归约。}
只须证明在其他假设下，\(H\) 此时为环面。由已经证明的
\SGARefLink{sga3:x:corollary:8-2}{8.2}，只须证明 \(H\) 的纤维为环面。
可假定 \(S\) 是完备 Noether 局部环 \(A\) 的谱。由 \SGARefLink{sga3:x:corollary:3:3}{3.3}，
对每个 \(n>0\)，存在在 \(S\) 上有限的乘法型群 \(T_n\) 及同构
\[
  (T_n)_s\simeq {}_n(H_s),
\]
其中 \(s\) 是 \(S\) 的闭点。如同 \SGARefLink{sga3:x:lemma:3:1}{3.1} 那样论证，
并利用 \(T_n\) 在 \(S\) 上有限，\footnote{\textit{编者注。}展开这一点。}
可知前述同构来自一个同态 \(u_n:T_n\to H\)，而且唯一确定。
（对 \(A\) 的 Artin 商取极限。）

此外，由唯一性，当 \(m\) 是 \(n\) 的倍数时，同态 \(u_n:T_n\to H\)
由 \(u_m:T_m\to H\) 限制到 \({}_n(T_m)\simeq T_n\) 得到。
由 \SGARefLink{sga3:ix:corollary:6:6}{IX 6.6}，各 \(u_n\) 为单态射，所以各 \(T_n\)
是 \(H\) 的子群，且当 \(m\) 是 \(n\) 的倍数时，有
\[
  T_n={}_n(T_m).
\]

因此，对每个 \(t\in S\)，各群 \((T_n)_t\) 是 \(H_t\) 的子群，型为
\((\mathbb{Z}/n\mathbb{Z})^r\)，其中 \(r=\dim H_s=\dim H_t\)，
且当 \(m\) 是 \(n\) 的倍数时，有
\[
  (T_n)_t={}_n(T_m)_t.
\]
于是，\(H_t\) 为环面由下列引理得到：

\par\medskip\noindent\SGAInlineRefTarget{sga3:x:lemma:8-6}\textbf{引理 8.6。---}\enspace
\emph{设 \(H\) 是代数闭域 \(k\) 上的平滑仿射代数群，\((T_n)_{n>0}\)
是 \(H\) 的乘法型子群族，使对每个整数 \(n>0\)，\(T_n\) 的型为
\((\mathbb{Z}/n\mathbb{Z})^r\)，且对 \(n\) 的每个倍数 \(m\)，有
\(T_n={}_n(T_m)\)。在这些条件下，\(H\) 含有一个维数至少为 \(r\) 的、包含各 \(T_n\)
的环面。因此，若 \(H\) 连通且维数至多为 \(r\)，则 \(H\) 是 \(r\) 维环面。}

这是关于仿射代数群的一道练习，将通过反复引用\emph{圣经〔Bible〕} 来处理。
只考虑 \(n\) 与特征互素的 \(T_n\)。设 \(K\) 为各 \(T_n\) 的并在 \(H\)
中的闭包，赋予诱导的约化结构。标准论证表明 \(K\) 是 \(H\) 的交换代数子群。
由\emph{圣经〔Bible〕}，\S4.5，定理~4，\(K\) 同构于积 \(K_u\times K_s\)，
其中 \(K_u\)“幺幂”，\(K_s\) 可对角化。\(K\) 的每个可对角化子群都包含于
\(K_s\)，故各 \(T_n\) 是 \(K_s\) 的子群，所以 \(K=K_s\)。写成 \(K=D(M)\)，
其中 \(M\) 是在 \(\mathbb{Z}\) 上有限型的通常交换群。包含 \(T_n\subseteq K\)
意味着 \(M\) 有一个同构于 \((\mathbb{Z}/n\mathbb{Z})^r\) 的商群。
由于这对每个与 \(k\) 的特征互素的整数 \(n\) 都成立（取某个固定素数的幂已经足够），
\(M\) 的秩至少为 \(r\)。因此 \(K\) 包含一个 \(r\) 维环面 \(T\)。
当 \(H\) 连通且维数为 \(r\) 时，得到 \(T=H\)，完成
\SGARefLink{sga3:x:lemma:8-6}{8.6} 的证明。现在\SGARefLink{sga3:x:theorem:8-1}{定理~8.1} 已证。

\par\medskip\noindent\SGAInlineRefTarget{sga3:x:remark:8-7}\textbf{注记 8.7。---}\enspace
使用 \SGARefLink{sga3:x:theorem:8-1}{8.1}，应不难给出
\SGARefLink{sga3:x:corollary:4-7}{4.7} 和 \SGARefLink{sga3:x:corollary:4-8}{4.8} 的类似推广。
更有趣的问题是，在 \SGARefLink{sga3:x:theorem:8-1}{8.1} 中去掉 \(H\) 的纤维仿射这一假设后研究其情形。
可以证明，这时存在 \(s\) 的开邻域 \(U\)，使 \(H|_U\) 交换，且对每个 \(t\in U\)，
几何纤维 \(H_t\) 是阿贝尔簇被环面扩张所得的群。
\footnote{\textit{编者注。}此处给出参考文献？}
当然，这类问题中可限于 \(S\) 为离散赋值环的谱、\(s\) 为其闭点、\(t\) 为其泛点的情形。

此结果可以如下推广。对代数闭域 \(k\) 上的每个平滑连通代数群 \(G\)，
Chevalley 的著名定理断言，\(G\) 唯一地是阿贝尔簇 \(A\) 被平滑连通仿射群
\(V\) 扩张所得的群。定义 \(G\) 的\emph{阿贝尔秩}
（分别为\emph{约化秩}、\emph{幂零秩}及\emph{半单维数}），记为
\(\rho_{\mathrm{ab}}(G)\)（分别为 \(\rho_r(G)\)、\(\rho_n(G)\) 及 \(d_s(G)\)），
为 \(A\) 的维数（分别为 \(G\) 的极大环面的维数、\(G\) 的 Cartan 子群
\footnote{\textit{编者注。}回顾“Cartan 子群”的定义。}的维数，以及 \(G\)
或等价地 \(V\) 对其根取商的维数）；这些概念均见\emph{圣经〔Bible〕}。
还定义\emph{幺幂秩}为
\[
  \rho_u(G)=\rho_u(V)
  =\rho_n(G)-\rho_r(G)-\rho_{\mathrm{ab}}(G).
\]
当基域不代数闭时，同样的名称和记号表示 \(G_{\bar{k}}\) 的相应不变量，
其中 \(\bar{k}\) 是 \(k\) 的代数闭包。

现在，设 \(G\) 是离散赋值环的谱 \(S\) 上的平滑群概形。
以 \(\rho_{\mathrm{ab}}\) 等（分别以 \(\rho'_{\mathrm{ab}}\) 等）表示
特殊纤维（分别为泛纤维）的不变量。则有不等式
\[
  \left\{
  \begin{aligned}
    \rho_{\mathrm{ab}}&\leq \rho'_{\mathrm{ab}},\\
    \rho_r+\rho_{\mathrm{ab}}&\leq
      \rho'_r+\rho'_{\mathrm{ab}},\\
    d_s&\leq d'_s,
  \end{aligned}
  \right.
  \qquad
  \left\{
  \begin{aligned}
    \rho_n&\geq \rho'_n,\\
    \rho_u&\geq \rho'_u.
  \end{aligned}
  \right.
\]
等价地，若 \(G\) 在任意基 \(S\) 上平滑且有限型，则函数
\[
  s\longmapsto \rho_{\mathrm{ab}}(s),\qquad
  s\longmapsto \rho_{\mathrm{ab}}(s)+\rho_r(s),\qquad
  s\longmapsto d_s(s)
\]
是下半连续的，而函数
\[
  s\longmapsto \rho_n(s),\qquad
  s\longmapsto \rho_u(s)
\]
是上半连续的。\footnote{\textit{编者注。}为这些结果给出参考文献？}

同样的结果很可能仍在不假定 \(G\) 在 \(S\) 上平滑、只假定它在 \(S\) 上平坦且有限呈示时成立，
只须对代数闭域 \(k\) 上的代数群 \(G\)，将 \(\rho_{\mathrm{ab}}(G)\) 等理解为
\(G_{\mathrm{red}}^0\) 的相应不变量。

本讨论班将在 \(G\) 在 \(S\) 上仿射，或更一般地具有仿射纤维时给出一些这类结果。
在此情形中，将验证 \(\rho_r\)、\(\rho_n\) 从而 \(\rho_u=\rho_n-\rho_r\)
的半连续性，以及在纤维为约化群的点 \(s\) 的邻域中 \(\rho_s\) 的连续性。\footnotemark[71]

若已经假定 \(G\) 交换，则可如下推广 \SGARefLink{sga3:x:corollary:8-2}{8.2}。

\par\medskip\noindent\SGAInlineRefTarget{sga3:x:theorem:8-8}\textbf{定理 8.8。---}\enspace
\footnote{\textit{编者注。}G.~Prasad 和 J.-K.~Yu 将此结果推广为不要求 \(G\) 交换，
而在假设~(a) 及结论中将“环面”换成“约化群”（附加某些其他假设）；
参见\ \([{\rm PY06}]\)，定理~6.2。}%
\emph{设 \(G\) 是交换 \(S\)-群预概形，在 \(S\) 上平坦且有限呈示，具有仿射连通纤维。
设 \(s\in S\)，并假定：
\begin{enumerate}
\renewcommand{\labelenumi}{\alph{enumi})}
\item 若 \(\bar{k}\) 表示 \(\kappa(s)\) 的代数闭包，则 \((G_{\bar{k}})_{\mathrm{red}}\) 是环面；
\item 存在 \(s\) 的泛化 \(t\)，使 \(G_t\) 在 \(\kappa(t)\) 上平滑。
\end{enumerate}
在这些条件下，存在 \(s\) 的开邻域 \(U\)，使 \(G|_U\) 为环面。}

\emph{（还应注意，若仅假定对 \(s\) 的每个泛化 \(t\)，\(G_t\) 仿射，
分别为连通，则容易推出在 \(s\) 的某个开邻域中的 \(t\) 也具有相同性状。）}

\medskip
\noindent\emph{\SGARefLink{sga3:x:theorem:8-8}{8.8} 的证明。}
只须证明 \(G_s\) 在 \(\kappa(s)\) 上平滑。
事实上，由于 \(G\) 在 \(S\) 上平坦且有限呈示，这将推出 \(G\) 在 \(s\)
的某邻域上方在 \(S\) 上平滑（参见\ \SGARefLink{sga3:x:lemma:3:5}{3.5}），
于是落入 \SGARefLink{sga3:x:theorem:8-1}{8.1} 的情形。

为证明 \(G_s\) 在 \(\kappa(s)\) 上平滑，通常步骤将问题归约到 \(S\) 仿射且 Noether 的情形。
选取态射 \(S'\to S\)，其中 \(S'\) 是离散赋值环的谱，将闭点映到 \(s\)，
将泛点 \(\eta\) 映到 \(t\)。这样归约到 \(S\) 本身为离散赋值环的谱的情形；
还可假定它完备，剩余域代数闭，且 \(s\)、\(t=\eta\) 分别为 \(S\) 的闭点与泛点。

因此，\(G\) 在 \(S\) 上平坦、分离且有限型；泛纤维 \(G_\eta\) 平滑且连通，
特殊纤维 \(G_0\) 满足
\[
  T_0=(G_0)_{\mathrm{red}}
\]
是环面。设 \(m>0\) 是与 \(s\) 处的剩余特征，从而也与 \(\eta\)
处的剩余特征互素的整数。则 \(m\cdot\operatorname{id}_G\) 在每个纤维上平展
（VII\textsubscript{A}，8.4），由于 \(G\) 在 \(S\) 上平坦，因而它平展
（SGA~1，I 5.9）。所以其核 \({}_mG\) 在 \(S\) 上平展；由于 \(G\) 分离
\begingroup
\renewcommand{\thefootnote}{\(\ast\)}
\footnote{由 Raynaud 定理，VI\textsubscript{B}，5.3。}
\endgroup
\addtocounter{footnote}{-1}%
\footnote{\textit{编者注。}见 8.5 中的编者注 (62)。}
且在 \(S\) 上有限型，\({}_mG\) 也如此。由于
\[
  ({}_mG)_0={}_m(T_0),
\]
其次数为 \(m^r\)，其中
\[
  r=\dim T_0=\dim G_0=\dim G.
\]
由此可知 \(({}_mG)_\eta={}_m(G_\eta)\) 的秩至少为 \(m^r\)，
利用 \SGARefLink{sga3:x:lemma:8-6}{8.6}，这已证明 \(G_\eta\) 是 \(r\) 维环面，
因为 \({}_mG\) 的两个纤维具有相同的秩，进而如同
\SGARefLink{sga3:x:lemma:8-5}{8.5} 一样，\({}_mG\) 在 \(S\) 上有限。
\footnote{\textit{编者注。}修改前句。}

由于 \(S\) 完备且剩余域代数闭，有限平展覆盖 \({}_mG\) 完全分裂。
因此，经过 \({}_mG_0\) 的每一点都有一个 \(G\) 在 \(S\) 上的截面；
特别地，\(G_0\) 中有 \(G\) 在 \(S\) 上的截面经过的点处处稠密。
现在使用下列引理。

\par\medskip\noindent\SGAInlineRefTarget{sga3:x:lemma:8-9}\textbf{引理 8.9。---}\enspace
\emph{设 \(S\) 是维数为 \(1\) 的局部 Noether 正则预概形，\(G\) 是平坦且局部有限型的
\(S\)-群预概形，且对 \(S\) 的每个极大点 \(\eta\)，\(G_\eta\) 在
\(\kappa(\eta)\) 上平滑（这蕴含 \(G\) 约化）。假定 \(G\) 的正规化 \(X\)
在 \(G\) 上有限（由 Nagata，当 \(S\) 是完备离散赋值环的谱时成立；参见\
EGA~IV\textsubscript{2}，7.7.4），令 \(G'\) 为 \(X\) 的开子集，由 \(X\) 在 \(S\) 上平滑的点组成。
采用这些记号，有：
\begin{enumerate}
\renewcommand{\labelenumi}{\alph{enumi})}
\item 若投影 \(G'\to S\) 满射，则 \(G'\) 携带唯一的 \(S\)-群预概形结构，
  使典范态射 \(G'\to G\) 为群同态。
\item 假定对 \(S\) 的每个闭点 \(s\)，\(G_s^0\) 中有平展拟截面经过的点集
  对于 Zariski 拓扑在 \(G_s^0\) 中稠密。则 \(X\) 正则，\textup{(a)} 的假设成立，且映射
  \[
    \Gamma(G'/S)\longrightarrow\Gamma(G/S)
  \]
  是双射。
\end{enumerate}}

在我们所关心的情形，此引理适用，并给出群同态 \(u:G'\to G\)，其中
\(G'\) 在 \(S\) 上平滑，而
\[
  u_\eta:G'_\eta\xrightarrow{\sim}G_\eta
\]
是同构。必要时将 \(G'\) 换成一个开子群后，可假定 \(G'_0\) 连通。
由于 \(u_0:G'_0\to G_0\) 诱导满射态射 \(G'_0\to T_0\)，而 \(T_0\)
是与 \(G'_0\) 同维的环面，容易推出 \(G'_0\) 是环面
（例如使用 \SGARefLink{sga3:x:lemma:8-6}{8.6}，或参阅\emph{圣经〔Bible〕}，\S7.3，定理~3(a)）。
因此由 \SGARefLink{sga3:x:corollary:8-2}{8.2}，\(G'\) 是环面。
由 \SGARefLink{sga3:ix:theorem:6:8}{IX 6.8}，\(\ker u\) 是 \(G'\) 的乘法型子群；
由于其泛纤维为单位群，它本身为单位群，故 \(u\) 为单态射。
由 \SGARefLink{sga3:viii:proposition:7:9}{VIII 7.9} 可知 \(u\) 为浸入。
由于它满射且 \(G\) 约化，\(u\) 为同构，完成了
\SGARefLink{sga3:x:theorem:8-8}{8.8} 的证明。

只须证明 \SGARefLink{sga3:x:lemma:8-9}{8.9}。
对于 \SGARefLink{sga3:viii:proposition:7:9:item:a}{(a)}，使 \(u:G'\to G\) 成为群同态的
\(G'\) 上群律的唯一性显然，因为 \(G'\) 泛纤维上的群律已知（可假定 \(S\) 不可约）。
为证明存在性，容易归约到 \(S\) 为局部，即某离散赋值环 \(A\) 的谱的情形。
由唯一性，且因整闭包与基的平展扩张交换，可以在 \(S\) 上作平展基扩张。
这将问题归约到 \(A\)“严格局部”的情形，即 Hensel 且剩余域可分闭。
同样的归约适用于 (b)；在其假设下，现在可以将“平展拟截面”换成“截面”。

由于 \(G'\) 在 \(S\) 上平滑，且 \(X\) 正规，预概形 \(G'\times_S X\)
正规（因为它在正规预概形 \(X\) 上平滑）。所以复合态射
\[
  G'\times_S X\longrightarrow G\times_S G\longrightarrow G
\]
分解为
\[
  p:G'\times_S X\longrightarrow X.
\]
将证明此态射限制到 \(G'\times_S X\) 的开子集 \(G'\times_S G'\) 时，诱导态射
\[
  G'\times_S G'\longrightarrow G'.
\]

因此只须证明 \(G'_0\times_k G'_0\) 映入 \(X_0\) 的开子集 \(G'_0\)。
只要验证对于 \(G'_0\) 的每个 \(k\)-值点 \(g'_0\)，态射
\[
  \tag{\(\times\)}\label{sga3:x:lemma:8-9:equation:times}
  h'_0\longmapsto p_0(g'_0,h'_0)
\]
从 \(G'_0\) 到 \(X_0\) 取值于 \(G'_0\)。由于 \(G'\) 在 \(S\) 上平滑，
且 \(S\) 为 Hensel，每个这样的 \(g'_0\) 都来自 \(G'\) 在 \(S\) 上的截面 \(g'\)。
于是态射 \((\times)\) 由态射
\[
  h'\longmapsto p(g',h')
\]
从 \(G'\) 到 \(X\) 诱导；后者本身由 \(G\) 的自同构 \(h\mapsto g\cdot h\)
通过结构搬运得到，即沿 \(G\) 的截面 \(g\) 的左平移，其中该截面是 \(g'\) 的像。
因此 \(h'\mapsto p(g',h')\) 本身是 \(X\) 的自同构，故将 \(G'\) 映入 \(G'\)。断言得证。

只须证明在 \(G'\) 上如此得到的合成律是群律。
结合律立即来自泛纤维上的结合律（该泛纤维同构于 \(G\) 的泛纤维）。
此外，\(S\)-预概形 \(G\) 的取逆自同构诱导 \(X\) 的自同构，因此保持 \(G'\)，
并诱导 \(G'\) 的自同构 \(\sigma\)。可以验证 \(\sigma\) 对于 \(G'\) 上的合成律
具有逆元运算的性质。事实上，这仍归结为验证涉及 \(G'\) 在 \(S\) 上的纤维幂的
某些图交换；由于它们在 \(S\) 上平滑，只须在泛纤维上核查交换性，那里显然成立。
这证明了 8.9 的 (a)。

下面证明 (b)。令 \(Z'\) 为使 \(\mathcal O_{X,x}\) 非正则的 \(x\in X\) 所成的集合。
由 Nagata 的定理（EGA~IV\textsubscript{2}，6.12.6），它是闭的，且显然包含于 \(X_0\)。
令 \(Z\) 为它在 \(G\) 中的像，因此是 \(G_0\) 的闭子集。
则 \(Z\) 在 \(G_0\) 中无处稠密，即不包含 \(G_0\) 的任何极大点 \(y\)。
事实上，\(G_0\) 在 \(G\) 中由方程 \(t=0\) 定义，其中 \(t\) 是离散赋值环
\(A\) 的素元。由主理想定理，\(\mathcal O_{G,y}\) 的维数为 \(1\)。
因此，对位于 \(y\) 上方的每个 \(x\in X\)，\(\mathcal O_{X,x}\) 的维数为
\(1\)，因而是离散赋值环，因为 \(X\) 正规，从而在余维 \(1\) 处正则。

另一方面，对 \(G\) 在 \(S\) 上的每个截面 \(g\)，集合 \(Z'\) 显然在由
\(G\) 上沿 \(g\) 的左平移通过结构搬运得到的 \(X\) 的自同构下稳定。
所以 \(Z\) 在由 \(g_0\) 定义的 \(G_0\) 上左平移下稳定。由假设，这些 \(g_0\)
的集合在 \(G_0^0\) 中稠密。由于 \(Z\) 在沿这些 \(g_0\) 的平移下稳定，且是无处稠密的闭子集，
立即推出 \(Z=\varnothing\)，从而 \(Z'=\varnothing\)，所以 \(X\) 正则。
因此，\(X\) 在有截面经过的每一点在 \(S\) 上平滑。
\(G\) 在 \(S\) 上的每个截面唯一地提升为 \(X\) 在 \(S\) 上的截面，
从而提升为 \(G'\) 在 \(S\) 上的截面。这特别适用于单位截面，证明了 \(G'\)
在 \(S\) 中的像是整个 \(S\)，即 (a) 的假设成立。
这完成了 \SGARefLink{sga3:x:lemma:8-9}{8.9}，从而也完成了
\SGARefLink{sga3:x:theorem:8-8}{8.8} 的证明。

\section*{\SGAInlineRefTarget{sga3:x:section:9}9. 增补}
\addcontentsline{toc}{section}{9. 增补}
\setcounter{section}{9}
\setcounter{subsection}{0}
\footnote{\textit{编者注。}此增补节由 Fabrice Orgogozo 撰写（依照 Ofer Gabber 的提示）。}

\subsection{单纯集、拓扑斯、群胚与拓扑空间}

\par\medskip\noindent\SGAInlineRefTarget{sga3:x:notation:9-1-1}\textbf{记号 9.1.1。---}\enspace
\emph{(1) 设 \(E\) 是单纯集。可将下列对象与之对应：
\begin{enumerate}
\renewcommand{\labelenumi}{(\arabic{enumi})}
\setcounter{enumi}{1}
\item 拓扑斯 \(T=E^{\sim}\)，由自然对应于各集合 \(E_i\) 的拓扑斯，
  通过 \([{\rm Del74}]\)，6.3.1 所述步骤得到
  （另见 \([{\rm Ill72}]\)，VI.5.2 和 SGA~4，VI.7）；
\item 群胚 \(G=\Pi(E)\)，其对象是 \(E_0\) 的元素（“顶点”），
  箭头定义见 \([{\rm GZ67}]\)，II.7；
\item 拓扑空间 \(X=|E|\)，这是称为“几何实现”（或“拓扑实现”）的胞腔复形
  （参见\ \emph{同处}，III.1）。
\end{enumerate}}

注意，\(T\) 上的层无非是 \(E\) 上的单纯集。

\subsection{局部常值层；下降数据}

\par\medskip\noindent\SGAInlineRefTarget{sga3:x:definition:9-2-1}\textbf{定义 9.2.1。---}\enspace
以下各类对象上的\emph{局部常值层}定义如下：
\begin{enumerate}
\renewcommand{\labelenumi}{(\arabic{enumi})}
\item 单纯集 \(E\) 上，是单纯集态射 \(E'\to E\)，使对每个 \(i\in\mathbb{N}\)
  及每个 \(e\in E_i\)，面算子 \(d\) 在纤维之间诱导同构
  \[
    E'_e\xrightarrow{\sim}E'_{d(e)}
  \]
  （参见\ \([{\rm AM69}]\)，\S10）；
\item 拓扑斯 \(T\) 上，是 \(T\) 的对象 \(F\)，使存在满态射 \(U\to 1\) 及同构
  \[
    F\times U\simeq f^*I\times U,
  \]
  其中 \(I\) 是集合，\(f:T\to\mathbf{Ens}\) 是终态射（参见\ SGA~4，IX.2）；
\item 群胚 \(G\) 上，是 \(G\) 上的预层，即从 \(G\) 到集合范畴的反变函子
  （参见\ \([{\rm GZ67}]\)，附录 I.1.2）；
\item 拓扑空间 \(X\) 上，是 \(X\) 上通常意义下局部常值的集合层。
\end{enumerate}
最后：
\begin{enumerate}
\renewcommand{\labelenumi}{(\arabic{enumi})}
\setcounter{enumi}{4}
\item 单纯集 \(E\) 上的\emph{下降数据}由 \(E_0^{\sim}\) 上的一个层 \(F\)
  （即 \(E_0\) 上的集合值函数）及同构
  \[
    \alpha:p_1^*F\xrightarrow{\sim}p_2^*F,
  \]
  组成，其中 \(p_1,p_2:E_1^{\sim}\to E_0^{\sim}\) 是由面诱导的态射，
  并要求满足通常的余循环关系（参见\ \SGARefLink{sga3:iv:definition:2:1}{报告~IV，2.1}%
  \begingroup
  \renewcommand{\thefootnote}{1}
  \footnote{PP：我将“FGA，\emph{下降技术〔Technique de descente〕} I，
  \SGARefLink{sga3:i:subsection:1-6}{1.6}”换成“\SGARefLink{sga3:iv:definition:2:1}{第 IV 篇，2.1}”。}
  \endgroup
  \addtocounter{footnote}{-1}%
  \ 及下文）。
\end{enumerate}

这五类对象之间的态射，以及相应的逆像函子，都以显然的方式定义。

\subsection{若干范畴等价}

设 \(E\) 是单纯集，\(T\)、\(G\)、\(X\) 分别为对应的拓扑斯、群胚与拓扑空间。

\par\medskip\noindent\SGAInlineRefTarget{sga3:x:proposition:9-3-1}\textbf{命题 9.3.1。---}\enspace
\emph{\(E\)、\(T\)、\(G\)、\(X\) 上的局部常值层范畴，以及 \(E\) 上的下降数据范畴，相互等价。}

\medskip
\noindent\emph{证明概要。}
用 \SGARefLink{sga3:x:proposition:9-3-1:item:1:03}{(1)} 至
\SGARefLink{sga3:x:proposition:9-3-1:item:5}{(5)} 表示前段定义的对象范畴。

\begin{itemize}
\item[\((1)\Longleftrightarrow(5)\).]
  这是 \([{\rm AM69}]\)，10.6 的特例（另见 \([{\rm GZ67}]\)，附录
  \SGARefLink{sga3:x:corollary:2:3}{I.2.3}，\([{\rm Fri82}]\)，?.5.6，或
  \([{\rm Ill72}]\)，VI.8.1.6）。将对象 \(E'\to E\) 对应于对 \((E'_0,\alpha)\)
  的函子给出范畴等价，其中 \(E'_0\) 视为 \(E_0\) 上的层，\(\alpha\) 是唯一的同构
  \[
    p_1^*E'_0\xrightarrow{\sim}p_2^*E'_0
  \]
  其在每个 \(y\in E_1\) 处的纤维由下列同构给出，其中其在两个投影下的像记为
  \(x_1\) 和 \(x_2\)：
  \[
    (E'_0)_{x_1}\xleftarrow{\sim}(E'_1)_y
      \xrightarrow{\sim}(E'_0)_{x_2}.
  \]
\item[\(\SGAInlineRefTarget{sga3:x:proposition:9-3-1:item:5}(5)\Longleftrightarrow(3)\).]
  这是显然的：与 \(E\) 对应的群胚 \(G\) 中定义态射的两个关系之一，就是余循环关系。
\item[\((1)\Longleftrightarrow(4)\).]
  参见\ \([{\rm GZ67}]\)，附录 I.3.2.1。
\item[\(\SGAInlineRefTarget{sga3:x:proposition:9-3-1:item:1:03}(1)\Longleftrightarrow(2)\).]
  须证明 \(E\) 上的一个对象在单纯意义下是局部常值层，当且仅当它作为
  \(T=E^{\sim}\) 上的层是局部常值层。这由以下事实得到：\(T\) 上的局部常值层
  \(F_\bullet\) 恰是笛卡尔层，即使箭头
  \[
    ([n]\to[m])^*F_m\longrightarrow F_n
  \]
  为同构的层。最后这个断言是关于单纯拓扑斯的一般事实的一个特例，
  再结合 \(E_0^{\sim}\) 上每个层都局部常值这一事实即可。
  （另见 \([{\rm Ill72}]\)，VI.8.1.6。）
\end{itemize}
\hfill 证毕。

对每个群 \(H\)，前述范畴等价诱导 \(H\)-挠子范畴之间的等价。
这里它们是配备 \(H\) 的作用的局部常值层，其纤维同构于 \(H\) 通过平移作用于自身。

\subsection{基本群与群胚}

由前述范畴等价，对于每个群 \(H\) 和每个单纯集 \(E\)，\(H\)-挠子的同构类集
\[
  H^1(E,H),\qquad
  H^1(E^{\sim},H),\qquad
  \pi_0\bigl(\operatorname{Hom}(\Pi(E),BH)\bigr),\qquad
  H^1(|E|,H)
\]
自然地一一对应。回顾，\(BH\) 表示与 \(H\) 对应的单对象群胚，
\(\pi_0\) 是将一个范畴对应于其对象同构类集的函子。

同样，若 \(e\) 是 \(E\) 的一点，前述等价诱导在 \(e\) 上平凡化的
\(H\)-挠子的同构类集之间的双射，分别记为
\[
  H^1(E\mathbin{\mathrm{rel}}e,H),\qquad
  H^1(E^{\sim}\mathbin{\mathrm{rel}}e^{\sim},H),\qquad
  \operatorname{Hom}\bigl(\pi_1(\Pi(E),e),H\bigr),\qquad
  H^1(|E|\mathbin{\mathrm{rel}}|e|,H).
\]
回顾，有
\[
  \pi_1(\Pi(E),e)=\operatorname{Isom}_{\Pi(E)}(e,e).
\]
当 \(H\) 变动时，在连通情形中，这些函子由一个记为 \(\pi_1(E,e)\) 的群表示。
群 \(\pi_1(E,e)\) 同构于 \(\pi_1(\Pi(E),e)\) 及 \(\pi_1(|E|,|e|)\)，
从而也同构于 Kan 所定义的单纯集的基本群
（例如参见\ \([{\rm May67}]\)，16.1，或 \([{\rm Ill72}]\)，I.2.1.1）。
还应回顾，\(H^1(E,H)\) 又同构于模共轭的到 \(H\) 的同态集
\(H^1(\pi_1(E,e),H)\)，也记为
\[
  \operatorname{Hom\,ext}(\pi_1(E,e),H).
\]

\subsection{锥}

\par\medskip\noindent\SGAInlineRefTarget{sga3:x:definition:9-5-1}\textbf{定义 9.5.1。---}\enspace
\emph{\SGAInlineRefTarget{sga3:x:definition:9-5-1:item:1}(1) 设 \(f:E'\to E\) 是单纯集态射。回顾
（例如参见\ \([{\rm Del74}]\)，6.3.1），\(C(f)\) 表示如下点化单纯集，
其次数 \(n\geq0\) 中的底集为}
\[
  E_n\amalg\left(\coprod_{i<n}E'_i\right)\amalg\star,
\]
\emph{其中 \(\star\) 是单点集。单纯映射的定义留给读者；次数至多为二的面在下文回顾。
（另见 \([{\rm GZ67}]\)，VI.2，其中给出一个点化的变体。）
\(C(f)\) 上的局部常值层范畴等价于 \(E\) 上配备其在 \(E'\) 上逆像的平凡化的局部常值层范畴。
单纯集 \(C(f)\) 由 \(\star\to C(f)\) 自然带有标记点。}

\emph{(2) 设 \(f:T'\to T\) 是拓扑斯态射。用 \(C(f)\) 表示对象为如下五元组的拓扑斯}
\[
  \bigl(F,F',A,\alpha:f^*F\to F',
                 \beta:A\to\Gamma(T',F')\bigr),
\]
\emph{其中 \(F\)（分别为 \(F'\)）是 \(T\)（分别为 \(T'\)）的对象，
\(A\) 是集合，\(\alpha\)（分别为 \(\beta\)）是 \(T'\)（分别为
\(\mathbf{Ens}\)）中的态射。（此构造的变体另见 \([{\rm Del80}]\)，4.3.4 及
\([{\rm Ill72}]\)，III.4。）\(C(f)\) 上局部常值层的范畴等价于 \(T\) 上
配备其在 \(T'\) 上逆像的平凡化的局部常值层范畴。
拓扑斯 \(C(f)\) 由将上述五元组映为集合 \(A\) 的纤维函子自然带有标记点。}

\emph{(3) 设 \(f:G'\to G\) 是群胚态射。用 \(C(f)\) 表示下图的余极限}
\[
  G\longleftarrow G'\longrightarrow B1,
\]
\emph{其中 \(B1\) 是单对象范畴（一个对象，一个箭头）。
\(C(f)\) 上的局部常值层范畴等价于 \(G\) 上配备其在 \(G'\) 上逆像的平凡化的局部常值层范畴。}

\emph{(4) 设 \(f:X'\to X\) 是拓扑空间态射。用 \(C(f)\) 表示下图的余极限}
\[
  \star\longleftarrow X'\times\{0\}
  \longrightarrow X'\times[0,1]
  \longleftarrow X'\times\{1\}
  \longrightarrow X.
\]
\emph{\(C(f)\) 上的局部常值层范畴等价于 \(X\) 上配备其在 \(X'\) 上逆像的
平凡化的局部常值层范畴。拓扑空间 \(C(f)\) 由 \(\star\to C(f)\) 自然带有标记点。}

\subsubsection*{9.5.2. 单纯映射的锥上的下降数据}

设 \(f:E'\to E\) 是单纯集态射，\(C(f)\) 是其锥
（参见上文 \SGARefLink{sga3:x:definition:9-5-1:item:1}{(1)}）。
用 \(p\)（分别用 \(q,r\)）表示 \(E'\)（分别为 \(E,C(f)\)）的面映射。
在陈述下面的命题前，先用 \(p\) 和 \(q\) 描述次数至多为二的面 \(r\)。
约定 \(p_{ij}\)（分别为 \(p_i\)）是对应于像为 \(\{i,j\}\)
（分别为 \(\{i\}\)）的递增映射的面映射：
\[
  E'_2=E'_{\{1,2,3\}}\longrightarrow
  E'_1=E'_{\{1,2\}}
  \quad
  \left(\text{分别为 }
  E'_1=E'_{\{1,2\}}\longrightarrow E'_0=E'_{\{1\}}\right)
\]
这里相应的递增映射为
\[
  \{1,2\}\longrightarrow\{1,2,3\}
  \quad
  \left(\text{分别为 }
  \{1\}\longrightarrow\{1,2\}\right)
\]
%
\begingroup
\renewcommand{\thefootnote}{2}
\footnote{此约定不同于目前的单纯约定，但更接近 Grothendieck 在下降理论中使用的记号。}
\endgroup
\addtocounter{footnote}{-1}%
对于 \(E\) 和 \(C(f)\) 的面，也采用同样约定。

态射
\[
  r_1:C(f)_1=E_1\amalg E'_0\amalg\star
  \longrightarrow C(f)_0=E_0\amalg\star
\]
（分别为 \(r_2\)）为：
\begin{itemize}
\item \(q_1:E_1\to E_0\) (分别为 \(q_2:E_1\to E_0\))，在 \(E_1\) 上；
\item \(E'_0\to\star\) (分别为 \(f_0:E'_0\to E_0\))，在 \(E'_0\) 上；
\item \(\star\to\star\)，在 \(\star\) 上。
\end{itemize}
同样，态射
\[
  r_{21}:C(f)_2=E_2\amalg E'_1\amalg E'_0\amalg\star
  \longrightarrow C(f)_1=E_1\amalg E'_0\amalg\star
\]
（分别为 \(r_{32},r_{31}\)）为：
\begin{itemize}
\item \(q_{21}:E_2\to E_1\) (分别为 \(q_{32},q_{31}\))，在 \(E_2\) 上；
\item \(p_1:E'_1\to E'_0\) (分别为
  \(f_1:E'_1\to E_1\), \(p_2:E'_1\to E'_0\))，在 \(E'_1\) 上；
\item \(E'_0\to\star\) (分别为
  \(\operatorname{Id}:E'_0\to E'_0\),
  \(\operatorname{Id}:E'_0\to E'_0\))，在 \(E'_0\) 上。
\end{itemize}

设 \(H\) 是 \(C(f)\) 上的下降数据，即
\[
  C(f)_0=E_0\amalg\star
\]
上的一个层 \(H_0\)，配备同构
\[
  \gamma:r_1^*H_0\xrightarrow{\sim}r_2^*H_0
\]
在其于下列集合上的逆像之间：
\[
  C(f)_1=E_1\amalg E'_0\amalg\star,
\]
并满足余循环关系
\[
  r_{31}^*\gamma=r_{32}^*\gamma\circ r_{21}^*\gamma.
\]

\(\gamma\) 在 \(E_1\) 上（分别在 \(E'_0\) 上）的限制是下降数据
\[
  \beta:q_1^*G_0\xrightarrow{\sim}q_2^*G_0,
\]
其中 \(G_0\) 是 \(H_0\) 在 \(E_0\) 上的限制（分别是平凡化
\[
  t:H_0(\star)\xrightarrow{\sim}f^*G_0=:F_0,
\]
其中 \(H_0(\star)\) 是茎为 \(H_0(\star)\) 的常层）。
\(\gamma\) 所满足的余循环关系限制到 \(E_2\)（分别到 \(E'_1\)）后，
是 \(\beta\) 的余循环关系（分别为
\[
  p_2^*(t)=f_1^*(\beta)\circ p_1^*(t)).
\]
后一关系表示，平凡化 \(t\) 与 \(\beta\) 在 \(F_0\) 上诱导的下降数据
\(f_1^*(\beta)\) 相容。由此可知：

\par\medskip\noindent\SGAInlineRefTarget{sga3:x:proposition:9-5-3}\textbf{命题 9.5.3。---}\enspace
\emph{设 \(f:E'\to E\) 是单纯集态射，\(C(f)\) 是其锥，以 \(\star\) 为基点。
\(C(f)\) 上在 \(\star\) 上平凡化的下降数据范畴，等价于 \(E\) 上配备其在
\(E'\) 上诱导的下降数据的平凡化的下降数据范畴。}

\subsection{函子的可表性：
\(\overline{\pi}^{\,1}(S'/S,\xi;{-})\)}

采用第 90 页的记号，从现在起假定
\[
  S_0=S'
  \qquad\text{和}\qquad
  S_1=S'\times_S S'
\]
的连通分量为开集。在此假设下，由前述回顾立即得到下列显式组合描述。

\par\medskip\noindent\SGAInlineRefTarget{sga3:x:proposition:9-6-1}\textbf{命题 9.6.1。---}\enspace
\emph{对每个群 \(G\)，带有 \(\xi\) 上平凡化的下降数据的同构类集
\(\overline{\pi}^{\,1}(S'/S,\xi;G)\) 自然地与相对非阿贝尔上同调集一一对应：}
\[
  H^1(K\mathbin{\mathrm{rel}}k,G)
  :=H^1(\widetilde K\mathbin{\mathrm{rel}}\star,G).
\]

因此，只要单纯集
\[
  \widetilde K=C(k\to K)
\]
连通，此函子就可由一个群表示，记为 \(\pi_1(K,k)\)，称为\emph{相对基本群}。
立即可验证，这当且仅当映射 \(\pi_0(k)\to\pi_0(K)\) 满射。
特别地，只要单纯集 \(K\) 连通就足够了。正如正文所指出，
当概形 \(S\) 连通且态射 \(S'\to S\) 普遍淹没时，确实如此。

\subsection{\(k\) 的可缩性；\(k_\bullet\to K_\bullet\) 单射性的反例}

下列两个命题使编者注 (53) 更为精确。

\par\medskip\noindent\SGAInlineRefTarget{sga3:x:proposition:9-7-1}\textbf{命题 9.7.1。---}\enspace
\emph{设 \(\kappa\) 是代数闭域，\(X\) 是连通 \(\kappa\)-概形。
对每个整数 \(i\geq0\)，以 \(X_i\) 表示 \(X\) 在 \(\kappa\) 上的第
\((i+1)\) 个纤维幂。对每个整数 \(i\geq0\)，典范映射}
\[
  \pi_0(X_i)\longrightarrow\pi_0(X)^{i+1}
\]
\emph{是双射。特别地，单纯集 \(\pi_0(X_\bullet)\) 可缩。}

只须证明下列引理，它由 Künneth 公式通过取极限得到。[展开？]

\par\medskip\noindent\SGAInlineRefTarget{sga3:x:lemma:9-7-2}\textbf{引理 9.7.2。---}\enspace
\emph{设 \(\kappa\) 是代数闭域，\(X\) 和 \(Y\) 是两个连通 \(\kappa\)-概形。
则纤维积 \(X\times_\kappa Y\) 连通。}

\par\medskip\noindent\SGAInlineRefTarget{sga3:x:proposition:9-7-3}\textbf{命题 9.7.3。---}\enspace
\emph{设 \(X\) 是连通概形，\(X'\to X\) 是有限平展态射，
\(\bar{x}\) 是 \(X\) 的一个位于点 \(x\) 的几何点。典范态射}
\[
  \pi_0(X'_{\bar{x}})\longrightarrow\pi_0(X')
\]
\emph{是满射。}

事实上，\(X'\) 的每个连通分量的像在 \(X\) 中既开又闭，
因而与纤维 \(X'_x\) 相交。因此
\[
  \pi_0(X'_x)\longrightarrow\pi_0(X')
\]
是满射，下列映射也如此：
\[
  \pi_0(X'_{\bar{x}})\longrightarrow\pi_0(X'_x).
\]

\section*{参考文献}
\addcontentsline{toc}{section}{参考文献}
\footnote{\textit{编者注。}本篇引用的补充文献；\([{\rm AM69}]\) 及以下各项涉及增补。}

\begin{description}
\item[\({[{\rm BAC}]}\)]
  N.~Bourbaki, \emph{交换代数〔Algèbre commutative〕}, 第~I--IV 章, Hermann, 1961;
  Masson, 1985; Springer-Verlag, 2006.
\item[\({[{\rm BEns}]}\)]
  N.~Bourbaki, \emph{集合论〔Théorie des ensembles〕}, 第~I--IV 章, Hermann, 1970.
\item[\({[{\rm Pes66}]}\)]
  C.~Peskine, \emph{Zariski“主定理”的一个推广〔Une généralisation du ``Main Theorem'' de Zariski〕},
  Bull.\ Sci.\ Math.\ \textbf{90} (1966), 119--127.
\item[\({[{\rm PY06}]}\)]
  G.~Prasad 和 J.-K.~Yu, \emph{关于拟约化群概形〔On quasi-reductive group schemes〕},
  J.\ Alg.\ Geom.\ \textbf{15} (2006), 507--549.
\item[\({[{\rm Ray70}]}\)]
  M.~Raynaud, \emph{Hensel 局部环〔Anneaux locaux henséliens〕}, Lect.\ Notes Maths.\
  \textbf{169}, Springer-Verlag, 1970.
\item[\({[{\rm AM69}]}\)]
  M.~Artin 和 B.~Mazur, \emph{平展同伦〔Etale homotopy〕}, Lect.\ Notes Maths.\
  \textbf{100}, Springer, 1969.
\item[\({[{\rm Del74}]}\)]
  P.~Deligne, \emph{Hodge 理论 III〔Théorie de Hodge III〕}, Publ.\ Math.\ IHÉS
  \textbf{44} (1974), 5--77.
\item[\({[{\rm Del80}]}\)]
  P.~Deligne, \emph{Weil 猜想 II〔La conjecture de Weil II〕}, Publ.\ Math.\ IHÉS
  \textbf{52} (1980), 137--252.
\item[\({[{\rm Fri82}]}\)]
  E.~M.~Friedlander, \emph{单纯概形的平展同伦〔Etale homotopy of simplicial schemes〕},
  Princeton Univ.\ Press, 1982.
\item[\({[{\rm GZ67}]}\)]
  P.~Gabriel 和 M.~Zisman, \emph{分式演算与同伦理论〔Calculus of fractions and homotopy theory〕}, Springer, 1967.
\item[\({[{\rm Ill72}]}\)]
  L.~Illusie, \emph{余切复形与形变 I \& II〔Complexe cotangent et déformations I \& II〕},
  Lect.\ Notes Maths.\ \textbf{239} \& \textbf{283}, Springer,
  1971 \& 1972.
\item[\({[{\rm Kan58}]}\)]
  D.~Kan, \emph{同伦群的一个组合定义〔A combinatorial definition of homotopy groups〕},
  Ann.\ of Math.\ \textbf{67} (1958), 282--312.
\item[\({[{\rm May67}]}\)]
  J.~P.~May, \emph{代数拓扑中的单纯对象〔Simplicial objects in algebraic topology〕},
  Univ.\ of Chicago Press, 1967.
\end{description}
